% =====================================================================================
% theory_core.tex -- theory core of the short draft of "Degrees of Separation" (author: Sora Yng).
% Written 2026-09-25; revised the same day after two referee reports (notes at the end of theory_appendix.tex).
% Condensed and restated from paper/main.tex Section 3, paper/si.tex Section S11 and THEORY_NOTES.md. Every number
% is copied from paper/main.tex, paper/si.tex, THEORY_NOTES.md or a repo-root *_RESULT.md; none is new. Proofs,
% the table of revealed rankings and the ledger of results: theory_appendix.tex (label app:thy).
% Needs amsmath, array, booktabs and natbib. Bibliography: ../refs_merged.bib plus refs_short.bib (entries there
% verified against Crossref). The environment and macro below are defined only if absent.
% Revised again 2026-09-25 for consistency with draft.tex (first-order SpringRank limit; T0 as a lower bound; T1;
% graduate study, coverage and location in the career rise; national-scale exceptions). draft.tex does not \input
% this file; it carries its own condensed Section 2.
% =====================================================================================
\makeatletter
\@ifundefined{thyprop}{\newtheorem{thyprop}{Proposition}}{}
\makeatother
\providecommand{\ci}[2]{$[#1,\allowbreak\,#2]$}

\section{One name, two markets}\label{sec:thy}

A university's name is valued in two markets: by the academic labour market, through the doctorates its
departments place (\emph{academic reputation}), and by the graduate labour market, through what employers pay its
bachelor's graduates (\emph{employer reputation}). We ask whether the two valuations agree
for a university $\times$ field programme and, for a study with admission-based variation, which one students follow
and which one predicts the causal return. \citet{anelli2020returns} estimates it at the admission
cutoff of a highly selective private university, partly through changed field choice, with degree completion,
on-time graduation and peer quality among the channels. Written after the descriptive
results, the model organises them and is tested only by checks T0--T8 (Appendix~\ref{app:thy}; specified before
the analysis; not publicly registered).

\paragraph{Primitives.} The unit is the programme $(i,f)$, over institutions $\mathcal I_f$ with a bachelor's
programme and a ranked doctoral department in field $f$. The two reputations value different
products sold under one name: doctorates and bachelor's graduates. Comparisons are across institutions
within a field: the field effect $\mu_f(t)$ is differenced out, so the model says nothing about the choice of field,
and $F$, a within-field rank, has no cross-field content. In deviations from field means,
\begin{equation}\label{eq:thy-latent}
s_{if}=\alpha_fa_i+r_{if},\qquad \theta_{if}=\beta_fx_i+y_{if},\qquad x_i=\xi_i+v_i,
\end{equation}
where $s$ is the academy's valuation of the department and $\theta$ the mean productivity of the programme's
graduates, selection included. University-wide are $a_i$ ($\mathrm{Var}\,a=1$), intake ability $\xi_i$ on the
dimension the admission criterion measures, and $v_i\perp\xi_i$, the part of graduate quality that measured intake
does not predict (unmeasured ability, value added, peer and resource effects);
$r_{if}$ and $y_{if}$ are field-specific, with $\kappa_f=\mathrm{Corr}_f(r,y)$; $\beta_f$ prices $x$. We assume
(A-orth) that $r$ and $y$ are uncorrelated with every university-wide variable.

\paragraph{Reputations.} \emph{Academic reputation} is the position revealed by the net direction of faculty hiring,
not a modelled belief about doctoral quality (Table~\ref{tab:thy-rankings}). In a gravity model of hires, SpringRank, the published $F$
\citep{wapman2022quantifying,debacco2018physical}, converges, to first order and up to scale, to $s^*=s+\pi/(2\beta)$, with $\pi=\log(p/h)$ its
production of doctorates relative to hiring and $\beta$ the hierarchy's steepness
(Proposition~\ref{prop:thy-springrank}). Directions identify only $s^*$; volume neutrality (H0: $\pi$ constant) is
the normalisation that reads $s^*$ as the valuation $s$. T0 finds $F$ partly production: coupling keeps $0.598$
\ci{0.471}{0.725} net of production, a lower bound, since residualising also removes net export that valuation
itself produces. T1 finds that $F$ does not order departments alike as placers and as hirers of PhDs (``not one
index''; its benchmark calibration was fixed after the literal one degenerated). \emph{Employer reputation} is the market's current valuation
$V_{if}(t)$, its mean belief about its graduates' productivity $t$ years out. Under
M1--M4 (Lemma~\ref{lem:thy-wage}) log median pay is
\begin{equation}\label{eq:thy-wage}
m_{if}(t)=\mu_f(t)+V_{if}(t)+\beta_fg_i(t),\qquad V_{if}(t)=(1-\lambda_t)R_{if}+\lambda_t\theta_{if},
\end{equation}
with entry reputation $R=V(0)$, weight $\lambda_t$ on graduates' own records and growth $g_i(0)=0$; the premium
decays, $V-\theta=(1-\lambda_t)(R-\theta)$ \citep{farber1996learning}. A name's record pools many graduates, so where
employers see them often $R\approx\theta$ at the name (Lemma~\ref{lem:thy-umbrella}): pay there ranks
selection-inclusive expected productivity, and belief and mean productivity cannot be separated. In the
closest economics usage, a college's reputation is its graduates' mean admission score \citep{macleod2017big}, our
third, student-revealed ranking; in that vocabulary the two valuations agree mostly through college reputation.

\paragraph{Students.} Student $k$ values programme $p$ at $\mathbb E_k[w_p(a_k)\mid\text{signals}]+\varepsilon_{kp}$
(signals such as rankings and programme earnings); admission is down to a cutoff, so cutoffs and
intake are equilibrium outcomes. Intake orders programmes as students do under a common ranking without idiosyncratic
tastes (S1; Lemma~\ref{lem:thy-students}), which the self-selection in \citet{kirkeboen2016field} contradicts,
and at the programme level only where admission is by programme, as in the UK; US students choose a major after entry, and our
US selectivity is university-level. Whether students follow academic reputation, employer reputation, peer ability
or value added is the next study's question, a university $\times$ field counterpart of
\citet{abdulkadiroglu2020parents}.

\begin{thyprop}[What agreement measures]\label{prop:thy-agree}
Under A0, classical errors ($F=s^*+e_F$, $Y=m+e_Y$, $S=\xi+e_S$), A-orth and Lemmas~\ref{lem:thy-umbrella}
and~\ref{lem:thy-wage}, the normal-score correlation of $F$ and $Y(t)$ is
$\sqrt{\mathrm{rel}_F\mathrm{rel}_Y}\,\mathrm{Corr}\{s^*,m(t)\}$, which the Spearman coupling $\rho_f(t)$ increases
with, and
\begin{equation}\label{eq:thy-decomp}
\mathrm{Cov}(s^*,m)=\underbrace{\frac{\mathrm{Cov}(s^*,\xi)\,\mathrm{Cov}(\xi,m)}{\mathrm{Var}\,\xi}}_{\text{through measured intake}}
+\underbrace{\alpha_f\mathrm{Cov}(a^{\perp\xi},m^U)}_{\text{university, beyond measured intake}}
+\underbrace{\Lambda_{ft}\kappa_f\sigma_r\sigma_y}_{\text{field-specific}}
+\underbrace{\frac{\mathrm{Cov}(\pi^{\perp\xi},m)}{2\beta}}_{\text{production}},
\end{equation}
with $m^U$ the university-wide part of pay and $\Lambda_{ft}=\omega_f+(1-\omega_f)\lambda_t$. The partial covariance
given $S$ is the last three terms plus $(1-\mathrm{rel}_S)$ times the first.
\end{thyprop}

The first term is not selection: peer quality and resources that rise with intake load on $\xi$
\citep{dale2002estimating}, and at the programme mean own and peer ability coincide. Selectivity is a common
component, a confounder if intake is fixed and a mediator if students sort on reputation
\citep{macleod2015reputation}. Most US agreement goes with institution-level
selectivity and composition: Pell share, control and the state earnings level take mean coupling from $+0.43$ to
$+0.31$, and the admission rate and average SAT to $+0.14$ (44 fields); the split depends on the order of entry, and
selectivity's Shapley share of $R^2$ is 0.173, against 0.076 for $F$. The UK controls graduates' own
prior attainment instead, whereas institution-mean SAT also removes peer and resource effects: coupling of $G$ keeps
66\% \ci{51}{79} among same-band graduates, 85\% \ci{80}{90} with band-standardised earnings, and 73\% \ci{67}{78}
of its log-earnings gradient ($+0.100$ per SD).

\begin{thyprop}[University or department]\label{prop:thy-level}
Let $c_X=\rho(X,Y)$ and $\cos\vartheta=\rho(F,G)$, and on standardised ranks let $\mathrm{sr}_F$ be the
correlation of $Y$ with the residual of $F$ on $G$. Then $c_F=c_G\cos\vartheta+\mathrm{sr}_F\sin\vartheta$, and $F$ out-predicts $G$ if
and only if $\mathrm{sr}_F>c_G\tan(\vartheta/2)$. If $G=a$ and H0 holds, $\mathrm{sr}_F$ prices only the
field-specific term, attenuated by a reliability below $\mathrm{rel}_F$.
\end{thyprop}

At the across-field means the threshold is $0.156$ against
$\mathrm{sr}_F=0.077$, exceeded in 15 of 57 fields. Within institutions, one SD of department prestige goes with
0.84\% higher four-year pay \ci{0.00}{1.64} (about zero with one- and five-year medians; up to about 3\% under the
lower-bound reliability), which need not be departmental, since $y_{if}$ also carries selection into the major
\citep{bleemer2022will}. We make no claim about why agreement below the name is weak.

\paragraph{At an admission cutoff.} Individually, the same model (Corollary~\ref{cor:thy-cutoff}) splits the effect
of admission to $p$ over the fallback $p'$ in the same field into
\begin{equation}\label{eq:thy-cutoff}
\tau(t)=\Delta\delta+(1-\lambda_t)(\Delta R-\Delta\delta)+\beta_f\Delta g(t),
\end{equation}
value added $\Delta\delta$ plus a reputation rent that learning erodes, at known names
$(1-\lambda_t)\beta_f\Delta\bar a$, the gap in mean ability, paid even if pure signal. Across fields add
$\Delta\mu_f(t)$; fallbacks differ across students, so the comparison is between programmes
\citep{hastings2013some,kirkeboen2016field}. Aggregates cannot test employer learning
(Proposition~\ref{prop:thy-time}); the time profile of $\tau(t)$ can. The design (Appendix~\ref{app:thy:design})
projects cutoff estimates of $\tau$ on target-minus-fallback gaps in the two reputations and peer ability.

\paragraph{Scope.} At known names, the rise of US coupling within fixed cohorts ($+0.031$ per year \ci{+0.021}{+0.041})
needs growth, falling pay variance unrelated to $F$, or learning about departments, small since $\mathrm{sr}_F$ is
(Proposition~\ref{prop:thy-time}); outside the model, graduate study, coverage and migration can produce the same
rise (with coverage and the research-doctorate rate as controls it is $+0.0103$ against $+0.0305$ per year, though
that rate tracks $G$; a fifth to a third moves with where graduates work, T7). Pay is not netted of location within
destination census divisions, and the Scorecard estimates are not netted of destination location at all.
Schedules that only rescale market pay leave correlations unchanged, so weak coupling under national scales (UK
nursing, medicine and teaching $+0.120$ against $+0.523$) comes from noisier medians, schedule variance unrelated to
status or weaker market agreement, not from the scaling; UK allied health ($+0.570$) and US special education
($+0.55$) couple strongly despite schedules. If fields differ only in the scale at which they price a common
university-level index, the field maps of coupling and selectivity pricing coincide
(Corollary~\ref{cor:thy-onefactor}; true-score $r=+0.94$, a consistency check). The model assumes, and the data do
not test, that pay equals expected productivity (M1); aggregates license no claim that attending raises pay or
that employers read rankings. Medians miss the upper tail where elite returns concentrate
\citep{chetty2023diversifying,zimmerman2019elite}, though net of selectivity and parental income $G$'s association
with the top 1\% share does not detectably exceed that with median pay (242 universities).
